%!TEX root = ../thesis.tex
\chapter{面向概念漂移的动态哈希检索方法}
\label{chap:scph}

\section{引言}
随着互联网多媒体内容的持续增长，海量图像数据的高效检索已成为信息处理与计算机视觉领域的重要研究问题。在实际应用中，基于关键词或人工标注的传统检索方式往往依赖完整、准确的标签或注释，而真实环境中仅有少量图像具备高质量的人工标注，这在很大程度上限制了其适用范围。面向该问题，基于内容的图像检索方法通过直接利用图像视觉特征实现相似性匹配，其中哈希检索因其在时间复杂度与存储开销方面的优势而受到广泛关注。哈希方法通常将高维特征空间中的样本映射到低维汉明空间，为每幅图像生成二值哈希码，并通过计算哈希码间的汉明距离实现近似最近邻检索；在理想情况下，语义或视觉内容相近的图像应具有相近的哈希码，从而在汉明空间中对应较小的距离。

然而，真实世界中的数据环境往往呈现显著的动态性：数据以流式或分批形式持续到达，既可能出现新的类别，也可能导致既有类别的数据分布随时间发生变化。上述随时间演化的分布变化可归纳为概念漂移现象，它会破坏“训练阶段获得的相似性结构在后续时间步保持稳定”的假设，使得许多针对静态数据环境设计的哈希模型难以长期维持稳定检索性能。为应对数据环境的动态变化，已有研究提出了多种动态哈希方法，通常通过在线或增量方式更新哈希模型以适应新到达数据。然而，这类方法在处理概念漂移时仍存在两个突出局限：其一，许多方法主要关注数据规模增长或数据动态到达这一事实，但对概念漂移这一更具挑战性的现象缺乏针对性机制；其二，当哈希模型在新时间步被更新后，往往需要访问历史数据并对数据库中已有样本的哈希码进行重新生成或整体更新。该过程在大规模数据库条件下代价高昂，不仅带来显著的计算与存储开销，也会削弱在线检索的实时性与系统效率。

在现有研究中，概念保留哈希（Concept Preserving Hashing, CPH）是能够同时面向概念漂移，并在更新哈希模型时无需更新旧数据哈希码的代表性方法。CPH 的核心思想是将每个类别对应到汉明空间中的固定区域，在数据环境变化时通过学习新的哈希函数将新到达数据投影至其所属类别的对应区域，使得历史数据可继续使用原有哈希码参与相似性评估，从而显著降低模型更新的系统开销。尽管 CPH 在检索效率方面表现突出，但其仍存在现实应用层面的限制：首先，CPH 属于全监督范式，要求新出现的数据具备完整标签，而在真实数据流中大量样本往往难以获得及时、准确的人工标注；其次，未标注数据并非“无信息”样本，其局部邻域结构蕴含可用于指导模型学习的潜在语义关系，忽略这部分信息会降低哈希学习的有效性；此外，CPH的性能在一定程度上依赖于对关键样本的精确识别与利用，这使其在复杂数据条件下可能面临稳定性挑战。

基于上述背景与问题，本章进一步引入一种基于概念保留的半监督动态哈希方法——Semi-supervised Concept Preserving Hashing（SCPH）。SCPH的基本思路是在保持“旧哈希码无需更新”这一高效更新机制的同时，将半监督学习引入动态哈希框架：一方面，针对未标注数据，SCPH依据其与附近标注样本的邻域关系为其分配伪标签（pseudo-label），从而在标注稀缺的条件下补充语义监督信号；另一方面，SCPH 为每个类别分配恒定的概念编码，用于记录类别在公共汉明空间中的相对位置，并以此为新到达数据构造目标编码与投影约束；在模型更新时，通过学习能够将新数据映射到汉明空间对应区域的新哈希函数，实现对新旧数据相似性关系的保持。由此，SCPH 旨在在更贴近真实应用的半监督与动态数据环境中，提升模型对概念漂移的适应能力，并在保证更新效率的同时维持较为稳定的检索性能。

\section{SCPH算法介绍}

SCPH在初始阶段，使用哈达玛矩阵（Hadamard matrix）为每个类别的数据生成概念编码。哈达玛矩阵是一种特殊的正交矩阵，其矩阵取值为 $+1$ 或 $-1$，且任意两行或两列之间两两正交。哈达玛矩阵的特殊性质使得以其为基础构建的概念编码在哈希比特层面具有较强的“去相关性”，从而符合哈希方法对于哈希编码均分性与判别性的基本期望，从而有利于在汉明空间中为不同类别提供稳定、可区分的概念参照。

具体而言，在初始阶段，生成一个哈达玛矩阵用于为数据中各个类别生成概念编码，该哈达玛矩阵以及各个类别的概念编码均保持固定且不再发生变化，从而为后续哈希函数的训练提供持续且稳定的指导。随后，SCPH利用未标注样本与其临近的标注样本之间的关系，为每个未标注样本分配伪标签。结合上述标签信息与各个类别的概念编码，为每个新到来的样本生成目标编码，并将所有的目标编码构成概念参考矩阵，用于描述新样本在汉明空间中的理想位置。在此基础上，通过保持新样本之间的相似性并将新数据投影至其对应的目标编码区域，学习得到新的哈希函数。

不失一般性的，我们将设数据块以相同规模按时间顺序逐批到达，并且在不同时间步长中，有标签数据样本和无标签数据样本的数量保持不变。在此设定下，在时间 $t = T$ 到达的数据块一批图像数据记为 $\mathbf{X} \in \mathbb{R}^{d \times n}$，其中 $d$ 表示每张图像的特征维度，$n$ 表示该数据块中图像的数量。对于任意时刻到达的数据块，有 $\mathbf{X} = [\mathbf{X}_l, \mathbf{X}_u]$，其中 $\mathbf{X}_l$ 表示有标签样本，$\mathbf{X}_u$ 表示无标签样本。在 $t$ 时刻，SCPH的目标是基于当前时刻新到来的数据块自适应的学习到一组新的哈希函数，即 $\mathbf{H}(\mathbf{x}) = \operatorname{sgn}(\mathbf{W}^T \mathbf{x})$。通过这组哈希函数，任意在时间 $t$ 到达的新图像样本 $\mathbf{x}_t$ 均可被在线映射成为一个 $r$ 位的二进制哈希编码 $\mathbf{b} \in \{-1, +1\}^r$，其中 $r$ 表示哈希编码的长度。哈希函数的具体形式如下：
\begin{equation}
  \mathbf{H}(\mathbf{X}) = \operatorname{sgn}(\mathbf{W}^T \mathbf{X})
\end{equation}
其中，$\mathbf{W} \in \mathbb{R}^{d \times r}$ 表示待学习的哈希投影矩阵，$\operatorname{sgn}(\cdot)$ 为符号函数。

\subsection{概念参考矩阵}

哈希方法通常期望语义相似的样本应具有相似的哈希码，从而在汉明空间中对应较小的汉明距离；反之，不相似样本的哈希码应尽可能区分开来。为提升编码的信息承载能力与判别性，哈希码各比特还应尽量相互独立，以避免冗余并提高码字的有效利用率。哈达玛矩阵在结构上能够较好地满足上述需求：其元素仅取 $-1$ 或 $+1$，且任意两行（或两列）两两正交，从而天然具备良好的去相关特性与码字分散性。因此，SCPH 首先引入哈达玛矩阵，并以其行向量为基础为各类别分配概念编码。对于 $2^k$ 阶哈达玛矩阵，其第 $i$ 行第 $j$ 列元素可表示为：
\begin{equation}
  M_{ij} = (-1)^{(i-1)(j-1)}
\end{equation}

进一步地，为使概念编码与哈希表示在维度上保持一致，SCPH 将哈达玛矩阵的阶数设定为大于等于哈希码长度 $r$ 且为 $2$ 的幂次，其定义如下：
\begin{equation}
  g = \min \{g \mid g=2^k, g \ge r, g \ge |L|, k=1,2,3,\dots \}
\end{equation}
其中，$|L|$ 表示当前已存在的类别数量。基于所构造的哈达玛矩阵，SCPH 从其行向量中以不放回随机抽取的方式为每一类别分配概念编码。具体而言，对某一类别随机选取一行作为该类别的概念编码，从而保证每个类别均对应唯一的概念编码。当有新类别出现时，则从哈达玛矩阵剩余未被选取的行中再随机抽取一行作为该新类别的概念编码。由于不同类别对应的行向量不同，各类别的概念编码彼此不相同。上述概念编码用于刻画各类别在汉明空间中的相对位置，为后续哈希函数的训练提供稳定的参照与指导。

在SCPH中，根据样本的标签为每个样本分配相应的目标编码。对于任一样本，其目标编码即为其所属类别对应的概念编码。对于无标签样本，则采用K近邻算法进行标签预测，为其分配伪标签。在 $t$ 时刻，新到来的有标签数据与无标签数据的欧式距离计算公式如下：
\begin{equation}
  d(\mathbf{x}_u, \mathbf{x}_l) = \|\mathbf{x}_u - \mathbf{x}_l\|_F
\end{equation}
其中，$\mathbf{x}_u$ 和 $\mathbf{x}_l$ 表示无标签样本和有标签样本，$\|\cdot\|_F$ 表示 Frobenius 范数。

对每个无标签样本，首先在当前新到达的标注样本中找到与其距离最近的 $K$ 个近邻。随后采用多数投票策略，将得票数最多的类别作为该无标签样本的伪标签（pseudo-label）。该无标签样本的伪标签可表示为：
\begin{equation}
  y_u = \arg\max_{c \in \{1, \dots, |L|\}} \sum_{i=1}^K \Psi(y_i, c)
\end{equation}
其中，$y_u$ 和 $y_i$ 分别表示未标注样本的伪标签以及第 $i$ 个近邻标注样本的真实标签；$c$ 表示类别标识。当第 $i$ 个近邻标注样本属于类别 $c$ 时，$\Psi(y_i, c) = 1$；否则 $\Psi(y_i, c) = 0$。在获得所有新到达样本的标签（包含伪标签）后，依据各类别的概念编码为这些样本生成目标编码，并将所有新样本的目标编码组合形成概念参考矩阵 $\mathbf{C} \in \{-1, +1\}^{n \times r}$。该矩阵刻画了每个新样本在汉明空间中的理想位置。

\subsection{SCPH的目标函数}

在 $t$ 时刻，SCPH的目标是在不更新旧数据的前提下，基于新到来的数据学习一组新的哈希函数，该哈希函数既能够为新到达的数据生成哈希编码，又能在概念漂移现象存在时，仍保留新数据与旧数据之间的相似性关系。因此，SCPH在学习新的哈希函数时，其目标函数综合考虑三个方面，即同构相似性、目标编码一致性以及哈希编码均分性。为简化表述，变量中的时间上标 $t$ 将被省略。SCPH学习新的哈希函数时，需要最大化的目标函数如下：
\begin{equation}
  \Theta(\mathbf{W}) = J_1(\mathbf{W}) + \alpha J_2(\mathbf{W}) + \beta J_3^*(\mathbf{W})
\end{equation}
其中，$\alpha$ 和 $\beta$ 均为非负超参数，$J_1(\mathbf{W})$ 表示同构相似性，$J_2(\mathbf{W})$ 表示目标编码一致性，$J_3^*(\mathbf{W})$ 表示哈希编码均分性。

在SCPH中，当两幅图像具有相同的标签（包含真实标签和伪标签）时，则这两幅图像语义相似；反之，则不相似。因此，新到来的数据的成对相似性矩阵 $\mathbf{S}$ 定义如下：
\begin{equation}
  S_{ij} = \begin{cases} +1, & \mathbf{x}_i \text{ 和 } \mathbf{x}_j \text{ 相似} \\ -1, & \mathbf{x}_i \text{ 和 } \mathbf{x}_j \text{ 不相似} \end{cases}
\end{equation}
矩阵 $\mathbf{S}$ 根据标签信息为依据对新数据块内的语义邻接结构进行编码，从而为后续同构相似项 $J_1(\mathbf{W})$ 的构造和优化提供基础。

对于新道来的数据块，我们期望哈希函数能够学习到新数据块内数据之间的相似性。具体来说，即相似的数据具有相似的哈希编码，并且在汉明空间中具有较小的汉明距离。所以，同构相似性可以通过如下式子进行量化评估：
\begin{equation}
  J_1(\mathbf{W}) = \operatorname{tr}(\mathbf{W}^T \mathbf{X} \mathbf{S} \mathbf{X}^T \mathbf{W})
\end{equation}

此外，由于动态数据环境中不可避免的存在概念漂移现象，同一类别的图像在特征空间中可能分布于不同位置。因此，SCPH在初始时刻为每个类别生成并固定了一个概念编码，使得每个类别在汉明空间中对应于一个固定区域。这些概念编码可以视为哈希函数在对新到来的数据进行映射时所应逼近的目标表征。因此，目标编码一致性项定义为：
\begin{equation}
  J_2(\mathbf{W}) = \operatorname{tr}(\mathbf{W}^T \mathbf{X} \mathbf{S} \mathbf{C})
\end{equation}
其中 $\mathbf{S}$ 和 $\mathbf{C}$ 分别表示新数据的成对相似性矩阵和概念参考矩阵。与此同时，理想的哈希编码还应该满足平衡性约束，即一串哈希编码中 $+1$ 与 $-1$ 的数量要尽可能的相等。哈希码的平衡性可以通过哈希编码的方差进行度量：
\begin{equation}
  J_3(\mathbf{W}) = \operatorname{tr}(\operatorname{sgn}(\mathbf{W}^T \mathbf{X}) \operatorname{sgn}(\mathbf{W}^T \mathbf{X})^T)
\end{equation}
由于 $\operatorname{sgn}$ 函数并不可导，对其进行松弛处理，得到：
\begin{equation}
  J_3^*(\mathbf{W}) = \operatorname{tr}(\mathbf{W}^T \mathbf{X} \mathbf{X}^T \mathbf{W})
\end{equation}

最终，SCPH的目标函数可表示为：
\begin{equation}
  \Theta(\mathbf{W}) = \operatorname{tr}(\mathbf{W}^T \mathbf{X} \mathbf{S} \mathbf{X}^T \mathbf{W}) + \alpha \operatorname{tr}(\mathbf{W}^T \mathbf{X} \mathbf{S} \mathbf{C}) + \beta \operatorname{tr}(\mathbf{W}^T \mathbf{X} \mathbf{X}^T \mathbf{W})
\end{equation}

通过令 $\Theta(\mathbf{W})$ 关于 $\mathbf{W}$ 的偏导为 $0$ 来求解最优解 $\mathbf{W}$。最优解 $\mathbf{W}$ 可通过如下式子求解得到：
\begin{equation}
  \mathbf{W} = \left( -\frac{\alpha}{2} \mathbf{X} \mathbf{S} \mathbf{X}^T + \beta \mathbf{X} \mathbf{X}^T \right)^{-1} \mathbf{X} \mathbf{S} \mathbf{C}
\end{equation}

在SCPH中，当哈希函数于后续时间步更新时，数据库中历史数据的哈希码无需随之重新计算或更新。在检索阶段，对于 $t$ 时刻的查询图像 $\mathbf{x}_q$，SCPH采用最新的 $K$ 组哈希函数为其生成 $K$ 个哈希码。随后，计算查询图像 $\mathbf{x}_q$ 的 $K$ 个哈希码与数据库中图像哈希码之间的累计汉明距离：
\begin{equation}
  D_{ham}(\mathbf{x}_q, \mathbf{x}_i) = \sum_{t=T-K+1}^T H_t(\mathbf{x}_q) \oplus B(\mathbf{x}_i)
\end{equation}
其中，$D_{ham}(\mathbf{x}_q, \mathbf{x}_i)$、$H_t(\cdot)$、$B(\cdot)$ 和 $\oplus$ 分别表示累计汉明距离、$t$ 时刻的哈希函数集合、待检索图像 $\mathbf{x}_i$ 的哈希码，以及异或运算。

\section{实验与分析}

\subsection{实验基本设置}

在本实验中，我们使用 CIFAR-10 和 CIFAR-100 两个数据集来模拟存在概念漂移现象的动态数据环境，以此验证 SCPH 算法的有效性。CIFAR-10 包含了 60,000 张 $32 \times 32$ 的彩色图像样本，样本被划分为 10 个语义类别；CIFAR-100 同样由 60,000 张 $32 \times 32$ 的 RGB 图像构成，包含 100 个细粒度类别，且可进一步归并为 20 个超类（super-classes），从而为更细粒度与层次化语义结构下的检索与哈希学习提供测试场景。CIFAR-10 和 CIFAR-100 均使用 512 维 GIST 特征进行表示。

在动态数据环境模拟中，本实验从时间 $t=0$ 至 $t=20$ 共生成 21 个数据批次，每个批次中，包含了 2000 张图像样本。其中，前 1000 张用于训练当前时刻的哈希函数，后 1000 张图像作为测试样本，用于评估当前时刻的检索性能。我们使用 CIFAR-10 数据集模拟存在新类出现的概念漂移的动态数据环境。具体而言，在 $t=0 \sim 5$ 时刻，每个数据批次从原始 10 个类别中随机抽取 5 个类别，并在这 5 个类别内随机采样 2000 个样本构成数据块。从 $t=6$ 时刻起，新的类别开始出现，新的数据批次的采样范围扩展至剩余 5 个未出现的类别，并与原先出现的类别共同参与随机采样，以形成后续时间步的新数据块，从而模拟数据流中新类别不断出现的场景。

对于数据分布发生改变的概念漂移场景，我们使用 CIFAR-100 进行动态数据场景模拟。具体而言，在 $t=0$ 时刻，我们先随机选择 10 个超类，然后从每个超类的第一个子类中随机抽取样本，一共抽取 2000 个样本；随后，在所选的 10 个超类中随机选择 6 个作为漂移超类；随着时间的推移，这些超类内部各个子类在数据批次中的采样比例会发生改变，其余 4 个超类的采样比例保持不变，以此来模拟数据分布发生改变的动态数据环境。具体的数据比例变化方式与 ICH 中的模拟设定一致。

在每个批次的 1000 个训练样本中，有标签数据的比例分别设置为 50\%、70\%、90\% 进行实验。SCPH 的两个超参数取值分别为 $\alpha=2$、$\beta=1$。此外，对于无标签样本的标签补全，本文采用 $K=20$ 的 KNN 算法，根据其在特征空间中的近邻关系为无标签样本分配伪标签。本文的哈希编码长度设置为 64 位。为降低随机性给实验结果带来的影响，本实验中所有实验均重复 10 次，以 10 次实验结果的平均值作为最终的实验结果。

在本实验中，SCPH 与静态哈希方法 EPH 和多种动态哈希方法进行对比，包括 OKH、OSH、HMOH、CIHR、CPH、OSelH、AQOH、UMH。对于静态哈希方法，其哈希函数仅在初始时刻 $t=0$ 训练得到，后续不再更新。因此，EPH 仅在 $t=0$ 时刻学习哈希函数，并在随后的时刻中直接使用初始时刻学习到的哈希函数将新到来的数据样本映射为哈希码，具体过程与原文一致。对于 OKH、HMOH、CPH、OSelH 等全监督方法，由于训练依赖训练样本的标签信息，因此在本实验中，这些方法仅利用数据块中有标签样本进行学习。对于无标签样本，则直接采用已学习到的哈希函数生成其哈希码。本实验的评估指标采用平均精度均值（Mean Average Precision, MAP）进行评估，并基于语义真值对不同方法的检索性能进行量化比较。

\subsection{实验结果分析}

在本实验中，我们模拟了新类出现和已有数据分布改变这两种概念漂移动态数据环境，并在每个批次中使用不同比例的有标签数据进行实验，以此来对比 SCPH 与其他方法在动态数据环境下的检索性能。

\subsubsection{新类出现}

图 3-1 至 图 3-3 展示了在新类出现的动态数据环境中，SCPH 在有标签数据占不同比例时，与其他方法在 $t=0 \sim 20$ 各个时刻的检索效果对比。首先，在 $t=0$ 时刻，SCPH 的检索效果并不是最好的，这是因为初始阶段，用于建模的监督信息和数据范围有限，哈希函数对于类别的结构信息学习得不够充分。随着时间的逐步推进，SCPH 在每个数据批次到来时都会学习新的哈希函数，使得模型能够学习到新数据中的数据分布信息，从而检索性能迅速提升。在 $t=5$ 时刻之后，数据分布发生概念漂移，新类别的持续引入使得检索性能在所有哈希方法上均出现下降。从理论上看，哈希学习通常依赖于历史数据所隐含的类别结构与相似性假设，而新类别的加入会破坏这一既有分布假设，进而改变汉明空间中的语义区域划分及样本间的相对相似关系。在此情况下，基于旧分布学习得到的哈希函数无法充分适配新的语义结构，导致新类样本在汉明空间中的嵌入产生偏移或混叠，从而显著降低整体检索性能。

\begin{figure}[htbp]
  \centering
  \includegraphics[width=0.8\textwidth]{scph/3-1}
  \caption{CIFAR-10 在每个批次中有标签数据比例为 50\%}
  \label{fig:cifar10-50}
\end{figure}

\begin{figure}[htbp]
  \centering
  \includegraphics[width=0.8\textwidth]{scph/3-2}
  \caption{CIFAR-10 在每个批次中有标签数据比例为 70\%}
  \label{fig:cifar10-70}
\end{figure}

\begin{figure}[htbp]
  \centering
  \includegraphics[width=0.8\textwidth]{scph/3-3}
  \caption{CIFAR-10 在每个批次中有标签数据比例为 90\%}
  \label{fig:cifar10-90}
\end{figure}

从实验结果可以观察到，随着时间推移及概念漂移的发生，SCPH 在各方法中始终保持最优的检索性能，表明其能够有效利用有标签与无标签样本信息以应对动态数据分布的变化。与同为半监督方法的 CIHR 相比，SCPH 在新类别出现后展现出更强的性能优势。相比之下，静态哈希方法 EPH 虽然在初始阶段具有较好的检索效果，但在动态数据环境中其性能随时间持续下降，且在概念漂移发生后下降更为显著。值得注意的是，在新类别引入后，SCPH 的检索性能虽出现短暂下降，但随后逐步回升，显示出其对概念漂移的良好适应能力。另一方面，全监督方法 CPH 在半监督数据设置下整体表现不佳，这主要源于其对有标签数据的强依赖，在标签数量有限时难以构建有效的概念参考矩阵。

\subsubsection{数据分布改变}

图 3-4 至图 3-6 展示了在数据分布随时间发生变化的动态数据环境下，SCPH 与多种对比方法在检索性能上的对比结果。总体来看，当数据分布发生改变时，SCPH 在各个时间阶段均表现出较为稳定且优于对比方法的检索效果，表明该方法在面对动态分布变化和概念漂移时具备较强的健壮性与自适应能力。

在每个数据块仅包含 50\% 标注样本的情况下，SCPH 初始 $t < 4$ 阶段时的检索性能并非最优的，这是由于半监督学习场景下，有标签的监督信息相对有限、模型尚未充分捕获数据分布变化规律。而随着时间的推进，新的数据样本不断到达，并且概念漂移现象发生了，数据的内在分布发生了改变，SCPH 的 MAP 值呈现不断上升的趋势，并最终超过所有的对比方法。相比之下，大多数基线方法在同一阶段表现出不同程度的性能下降，进而导致检索性能随时间的推移不断下降。

\begin{figure}[htbp]
  \centering
  \includegraphics[width=0.8\textwidth]{scph/3-4}
  \caption{CIFAR-100 在每个批次中有标签数据比例为 50\%}
  \label{fig:cifar100-50}
\end{figure}

\begin{figure}[htbp]
  \centering
  \includegraphics[width=0.8\textwidth]{scph/3-5}
  \caption{CIFAR-100 在每个批次中有标签数据比例为 70\%}
  \label{fig:cifar100-70}
\end{figure}

\begin{figure}[htbp]
  \centering
  \includegraphics[width=0.8\textwidth]{scph/3-6}
  \caption{CIFAR-100 在每个批次中有标签数据比例为 90\%}
  \label{fig:cifar100-90}
\end{figure}

从算法层面来看，SCPH 通过在哈希函数更新过程中引入基于“概念”的约束，对新旧数据之间的相似性关系进行统一建模，使哈希函数在随时间更新时保持相对稳定的映射行为，从而在一定程度上抑制了概念漂移引起的汉明空间结构变化。相较之下，缺乏此类约束的对比方法在数据分布持续变化时更容易出现哈希函数不稳定的问题，导致相似样本在汉明空间中的邻域关系逐渐失真，进而引发检索性能下降。

进一步地，随着有标签数据比例的逐步提高，SCPH 相较其他方法获得了更为显著的性能增益。一方面，更多的真实标注样本直接增强了监督信号的强度，有助于学习更加判别性的哈希函数；另一方面，充足的标注信息也提高了基于 KNN 的伪标签推断的可靠性，使未标注样本能够提供更准确且稳定的语义约束，从而促使模型学习到更加一致、信息量更高的哈希表示，最终带来整体检索性能的进一步提升。

\section{本章小结}

本章围绕存在概念漂移的动态数据环境下图像哈希检索所面临的关键问题，提出了一种半监督概念保留哈希方法（SCPH）。该方法在初始阶段为各类别分配唯一的概念编码，并利用 KNN 算法对未标注样本进行伪标签推断；随后，结合概念编码与真实/伪标签构建概念参考矩阵，从而在哈希函数学习过程中引入显式的概念约束，以增强模型在动态数据流场景下对语义变化的适应能力。大量实验结果表明，在存在概念漂移的动态数据环境下，SCPH 相较多种现有的哈希方法能够取得更为稳定且优越的检索性能，这验证了引入概念约束并充分利用半监督信息在缓解概念漂移、保持汉明空间结构稳定性方面的有效性。

尽管如此，SCPH 仍存在一定的局限性。该方法在每一个时间步均重新学习哈希函数，而实际应用中数据分布的变化程度可能较为有限，甚至在部分时间段内并未发生显著概念漂移，从而可能引入不必要的计算开销，影响其在线应用效率。针对上述问题，未来工作可在 SCPH 框架中引入概念漂移检测机制，以自适应地判定模型更新的触发时机与更新频率，在保证检索性能的同时进一步提升算法的计算效率与实用性。

需要指出的是，本章研究主要针对单模态图像数据，在动态数据环境下通过引入概念约束来提升哈希检索对分布变化的适应能力。然而，在实际应用场景中，数据往往以多模态形式存在，不同模态之间不仅同样面临语义一致性保持的问题，还进一步引入了跨模态语义对齐的挑战。因此，下一章将在本章研究思路的基础上，进一步拓展至多模态数据环境，研究基于多表结构的跨模态哈希检索方法，以实现不同模态间的有效语义关联与高效检索。
